\documentclass[11pt,a4paper]{article}
\usepackage[utf8]{inputenc}
\usepackage[T1]{fontenc}
\usepackage{amsmath,amssymb,amsthm,mathtools}
\usepackage{geometry}
\usepackage{booktabs}
\usepackage{hyperref}

\geometry{margin=1in}

\newtheorem{theorem}{Theorem}[section]
\newtheorem{corollary}[theorem]{Corollary}

\title{\textbf{Machine-Checked Symplecticity, Time-Reversibility, and Nine-Integral Conservation of the Pairwise Two-Body Decomposition for the Gravitational Three-Body Problem in Lean~4}}

\author{
  \textbf{Mohammad Hossam Mohammed Ali Zahran (\texttt{MHMZAHRAN})}\\
  Menoufia National University, Egypt\\
  \texttt{mhmzahran}
}

\date{November 2026}

\begin{document}

\maketitle

\begin{abstract}
Splitting the gravitational $N$-body problem into a composition of exact analytical Two-Body Kepler problems enables collisional orbital integration without coordinate singularities. While Hernandez and Bertschinger (2015) demonstrated numerically via Richardson-extrapolated Jacobian matrices that symmetric compositions of pairwise two-body Kepler propagators are symplectic on the full phase space, a machine-checked mathematical verification of this property and its associated conservation laws has remained absent from formal theorem-proving libraries. In this paper, we present the first complete formalization in the \textbf{Lean~4} theorem prover (paired with a zero-dependency C99 reference implementation) of the Three-Body Problem modeled as three interconnected Two-Body problems. All theorems are proved without \texttt{sorry} axioms.
\end{abstract}

\section{The Master Interconnected Two-Body System}

Let $\mathbf{r}_{12} = \mathbf{q}_1 - \mathbf{q}_2$, $\mathbf{r}_{23} = \mathbf{q}_2 - \mathbf{q}_3$, $\mathbf{r}_{31} = \mathbf{q}_3 - \mathbf{q}_1$, and $\kappa_{ij} = G / \|\mathbf{r}_{ij}\|^3$.

\begin{theorem}[Master Interconnected Two-Body Decomposition]
Let $M = m_1 + m_2 + m_3$. Subject to $\mathbf{r}_{12} + \mathbf{r}_{23} + \mathbf{r}_{31} = \mathbf{0}$, the relative acceleration of pair $(1,2)$ satisfies
\begin{equation}
  \ddot{\mathbf{r}}_{12} = -M \kappa_{12} \mathbf{r}_{12} + m_3 \Big[ (\kappa_{31} - \kappa_{12})\mathbf{r}_{31} + (\kappa_{23} - \kappa_{12})\mathbf{r}_{23} \Big].
\end{equation}
\end{theorem}

\section{Formal Verification of Global Symplecticity on \texorpdfstring{$\mathbb{R}^{18}$}{R18}}

\begin{theorem}[Canonicity of the Pairwise Kepler Embedding]
For any pair $(i,j)$ with normalized mass weights $\alpha + \beta = 1$, transforming $(\mathbf{q}_i, \mathbf{p}_i, \mathbf{q}_j, \mathbf{p}_j)$ into pair center-of-mass coordinates $(\mathbf{R}_{ij}, \mathbf{P}_{ij})$ and relative Kepler coordinates $(\mathbf{r}_{ij}, \mathbf{p}_{\mathrm{rel},ij})$ preserves the canonical symplectic $2$-form $\omega_{12}$ identically. Consequently, the $5$-stage Strang composition $\Psi_h = \Phi_{h/2}^{(12)} \circ \Phi_{h/2}^{(23)} \circ \Phi_h^{(31)} \circ \Phi_{h/2}^{(23)} \circ \Phi_{h/2}^{(12)}$ preserves $\omega_{18}$ identically, is strictly time-reversible, and conserves all $9$ Poincar\'e momentum, angular momentum, and barycentric boost invariants.
\end{theorem}

\end{document}
